\documentclass[conference]{IEEEtran}

\usepackage[utf8]{inputenc}
\usepackage[T1]{fontenc}
\usepackage[brazil]{babel}
\usepackage{graphicx}
\usepackage{booktabs}
\usepackage{multirow}
\usepackage{subcaption}
\usepackage{amsmath, amssymb}
\usepackage{float}
\usepackage{xcolor}
\usepackage{fancybox}

\graphicspath{{../assets/elections/plots/}{../assets/elections/tables/}}

% Placeholder para tabelas (substituir por screenshot)
\newcommand{\tablePlaceholder}[1]{%
  \begin{center}
    \fbox{\parbox{0.85\columnwidth}{%
      \centering\vspace{0.4cm}
      \textbf{[PLACEHOLDER]} \\[0.2cm]
      \textit{#1} \\[0.4cm]
    }}
  \end{center}
}

% -------------------------------------------------------------------------
\title{Análise Estatística Descritiva do Dataset de Eleições:\\
Financiamento de Campanha e Padrões Eleitorais dos Candidatos a\\
Deputado Federal nas Eleições Brasileiras de 2014}

\author{
  \IEEEauthorblockN{Heitor Teixeira - 2651254}
  \IEEEauthorblockA{
    \textit{MBA: Ciência de dados}\\
    UNIFOR\\
    Fortaleza, Brasil\\
    heitor.mba.unifor@gmail.com
  }
}

\begin{document}

\maketitle

% -------------------------------------------------------------------------
\begin{abstract}
Este trabalho apresenta uma análise estatística descritiva do dataset de
eleições brasileiras de 2014 (D2), com foco nos candidatos ao cargo de
Deputado Federal. O dataset integrado compreende 6.353 candidatos, com
dados de gastos de campanha, votos obtidos, resultado eleitoral e
características demográficas (gênero e raça), enriquecidos a partir de
registros oficiais do Tribunal Superior Eleitoral (TSE). A análise
inclui: classificação das variáveis; estatísticas descritivas completas
das variáveis quantitativas Gasto e Votos; distribuições de frequência;
análise de gasto por partido e por estado; estudo de fornecedores e setores
econômicos; coeficiente de correlação de Pearson entre gasto e votos em nível
nacional e estadual; e análise de representatividade por gênero e raça.
Os resultados revelam uma distribuição de gastos fortemente assimétrica à
direita (razão média/mediana de 22$\times$ em nível nacional), com o Amazonas
(AM) exibindo a correlação mais forte entre gasto e votos ($r = 0{,}902$),
seguido por Goiás (GO, $r = 0{,}605$) e Rio de Janeiro (RJ, $r = 0{,}556$).
A análise de equidade evidencia disparidades estruturais: candidatas mulheres
concentraram apenas 9{,}6\% dos gastos totais, correspondentes a 29{,}6\% dos
candidatos; candidatos brancos foram o único grupo com proporção de gastos
superior à sua representação (razão 1{,}40$\times$), enquanto candidatos
indígenas acessaram apenas 5\% do esperado.
\end{abstract}

\begin{IEEEkeywords}
Estatística descritiva, eleições brasileiras, financiamento de campanha,
correlação de Pearson, gênero, raça, análise exploratória de dados.
\end{IEEEkeywords}

% =========================================================================
\section{Introdução}

Este trabalho realiza uma análise de estatística descritiva sobre o dataset
de Eleições (D2), em atendimento à atividade avaliativa da disciplina de
Estatística Descritiva. O dataset reúne informações de candidatos ao cargo
de Deputado Federal nas eleições brasileiras de 2014, incluindo gastos de
campanha declarados, votos obtidos, resultado eleitoral e atributos
sociodemográficos provenientes do cadastro oficial do TSE.

A análise foi implementada em Python com as bibliotecas \textit{pandas},
\textit{matplotlib}, \textit{numpy}, \textit{scipy} e \textit{requests},
e cobre os seguintes requisitos: (i) classificação das variáveis qualitativas
e quantitativas; (ii) estatísticas descritivas completas de Gasto e Votos;
(iii) histogramas e distribuições de frequência; (iv) gráficos de barras para
variáveis qualitativas; (v) análise de fornecedores e setores econômicos;
(vi) coeficiente de correlação de Pearson entre Gasto e Votos em nível nacional
e para os estados RJ, GO e AM; e (vii) análise complementar de representatividade
e desempenho eleitoral por gênero e raça.

O formato de entrega adotado é o de artigo científico segundo o padrão IEEE,
pela mesma justificativa apresentada na análise do dataset D1: trata-se de uma
estrutura consolidada que favorece a organização clara dos resultados e sua
reprodutibilidade.

% =========================================================================
\section{Fundamentação Teórica}

\subsection{Estatística Descritiva}

A estatística descritiva visa sumarizar e apresentar sistematicamente as
características de um conjunto de dados, sem inferir parâmetros populacionais.
No contexto de dados eleitorais com alta variabilidade, a escolha das medidas
de posição e dispersão é particularmente relevante, pois distribuições
assimétricas tornam a média uma medida enganosa da tendência central típica.

\subsection{Medidas de Posição e Dispersão}

A média aritmética é sensível a valores extremos (\textit{outliers}), sendo
inadequada como única medida de posição em distribuições fortemente assimétricas.
A mediana, por ser resistente a valores extremos, representa com mais fidelidade
o candidato típico em distribuições de gastos de campanha. O desvio-padrão e o
coeficiente de variação complementam a análise ao quantificar a dispersão relativa
e absoluta dos dados.

\subsection{Distribuições de Frequência e Histogramas}

O histograma representa a distribuição de frequências de uma variável contínua
por retângulos adjacentes, com base no intervalo de classe e altura proporcional
à frequência. O número de classes foi definido pela regra da raiz quadrada
($k \approx \sqrt{n}$), adotada neste trabalho. A visualização em escala
logarítmica é necessária para variáveis com amplitude muito elevada, como gastos
de campanha, que variam de R\$\,0{,}72 a R\$\,8{,}4 milhões em nível nacional.

\subsection{Correlação de Pearson}

O coeficiente de correlação de Pearson $r$ mede a intensidade e a direção da
relação linear entre duas variáveis contínuas:
\begin{equation}
  r = \frac{\sum_{i=1}^{n}(x_i - \bar{x})(y_i - \bar{y})}
           {\sqrt{\sum_{i=1}^{n}(x_i - \bar{x})^2 \cdot
                  \sum_{i=1}^{n}(y_i - \bar{y})^2}}
\end{equation}
Valores de $|r|$ próximos de 1 indicam forte relação linear; valores próximos
de 0, ausência de relação linear. Neste trabalho, o teste de significância
associado avalia a hipótese nula $H_0: r = 0$ entre gasto de campanha e votos
obtidos.

\subsection{Análise de Representatividade}

A razão de representatividade é calculada como a proporção do total de gastos
de um grupo dividida pela proporção desse grupo no total de candidatos. Uma
razão igual a 1 indica distribuição proporcional; razão superior a 1 indica que
o grupo acessa mais recursos do que sua presença na corrida eleitoral; razão
inferior a 1 indica sub-representação financeira. Essa métrica permite comparar
grupos de diferentes tamanhos em condições de igualdade.

% =========================================================================
\section{Metodologia}

\subsection{Fontes de Dados e Integração}

O dataset principal foi fornecido pelo professor e contém registros de 6.353
candidatos a Deputado Federal nas eleições brasileiras de 2014, com variáveis
de gasto de campanha (Gasto), votos obtidos (Votos) e resultado eleitoral
(Situacao). Para enriquecer a análise demográfica, esses dados foram integrados
ao cadastro oficial de candidatos do TSE (2014), que fornece gênero, raça,
escolaridade e estado civil. Um segundo conjunto de dados com 366.289 registros
de transações com fornecedores foi utilizado para a análise setorial.

\subsection{Validação e Consistência dos Dados}

Os gastos do dataset do professor foram confrontados com os gastos oficiais do
TSE para aferir a confiabilidade da fonte. Das 5.319 candidaturas com dados
nas duas fontes, 95{,}02\% apresentaram correspondência exata (diferença
$<$ R\$\,0{,}01); apenas 4{,}96\% (264 candidatos) exibiram divergência
superior a R\$\,1{,}00, com diferença máxima de R\$\,966{,}926 e diferença
média de R\$\,585{,}09. Cerca de 40\% das divergências (107 casos) concentraram-se
no Rio de Janeiro, o que sugere que a base do professor pode ter sido compilada a
partir de uma versão anterior dos dados do TSE. As divergências identificadas não
comprometem as análises realizadas.

\subsection{Filtragem e Definição da Amostra}

Candidatos com Gasto igual a zero foram removidos das análises de distribuição
e correlação. Essa decisão baseia-se no fato de que nenhum candidato eleito
registrou gasto zero, e a inclusão desses registros distorceria as medidas de
tendência central e a correlação com votos. A amostra resultante para análise
de gasto compreende \textbf{5.323 candidatos} em nível nacional.

\subsection{Procedimentos Analíticos}

A análise foi organizada em seis etapas complementares: (i) classificação das
variáveis e caracterização inicial do dataset; (ii) cálculo de estatísticas
descritivas completas para Gasto e Votos; (iii) construção de histogramas e
análise da distribuição de gastos; (iv) análise de gastos por partido e por
estado; (v) análise de correlação de Pearson entre Gasto e Votos em nível
nacional e para os estados RJ, GO e AM; (vi) análise setorial de fornecedores;
e (vii) análise complementar de representatividade e desempenho eleitoral por
gênero e raça.

% =========================================================================
\section{Resultados e Discussão}

\subsection{Caracterização Inicial do Dataset}
\label{subsec:caracterizacao_inicial}

\begin{figure}[t]
  \centering
  \tablePlaceholder{Identificação e classificação das variáveis do dataset de Eleições.}
  \caption{Identificação e classificação das variáveis do dataset de Eleições 2014.}
  \label{tab:variaveis}
\end{figure}

A Tabela~\ref{tab:variaveis} apresenta a classificação das variáveis do dataset.
As variáveis quantitativas principais são Gasto (contínua, em R\$) e Votos
(discreta, contagem). As variáveis qualitativas incluem Situacao (3 categorias:
Eleito, Suplente, Não Eleito), Genero (2 categorias: Masculino, Feminino),
Raca (5 categorias: Branca, Parda, Preta, Amarela, Indígena), Partido e Estado.

\begin{figure}[t]
  \centering
  \includegraphics[width=\columnwidth]{table_2_3_1.png}
  \caption{Cruzamento de completude entre o dataset do professor e os dados
           oficiais do TSE 2014.}
  \label{tab:completude}
\end{figure}

A Tabela~\ref{tab:completude} evidencia que 84{,}87\% dos candidatos foram
localizados na base oficial do TSE. Os 15{,}06\% não localizados com Gasto = 0
não afetam as análises, pois o filtro de candidatos sem gasto registrado é
aplicado previamente.

\subsection{Estatísticas Descritivas}

A Tabela~\ref{tab:descritivas_br} apresenta as principais estatísticas
descritivas das variáveis quantitativas em nível nacional.

\begin{figure}[t]
  \centering
  \tablePlaceholder{Estatísticas descritivas de Gasto e Votos --- Brasil (6.353 candidatos).}
  \caption{Estatísticas descritivas das variáveis quantitativas --- Brasil.}
  \label{tab:descritivas_br}
\end{figure}

Os principais resultados são:

\begin{itemize}
  \item \textbf{Gasto}: média de R\$\,184.664 ($\sigma = $ R\$\,584.669),
        mediana de R\$\,5.804. A razão média/mediana de $\approx 22\times$
        evidencia forte assimetria positiva e indica que a mediana representa
        melhor o candidato típico.
  \item \textbf{Votos}: média de 18.058, mediana de 1.775; máximo de 1.524.361.
        Mesmo grau de assimetria, com desvio-padrão superando a média.
  \item \textbf{Percentil 75 do Gasto}: R\$\,46.569 --- 75\% dos candidatos
        gastaram menos de R\$\,47 mil, enquanto a média é puxada pelos poucos
        que gastaram milhões.
\end{itemize}

\subsection{Distribuição de Frequências: Gasto de Campanha}

A Figura~\ref{fig:hist_gasto} exibe o histograma dos gastos de campanha em
nível nacional (candidatos com Gasto $>$ 0), com eixo $x$ em escala logarítmica.
O número de classes foi determinado pela regra da raiz quadrada
($k \approx \sqrt{5323} \approx 72$).

\begin{figure}[t]
  \centering
  \includegraphics[width=\columnwidth]{distribuicao_gasto.png}
  \caption{Distribuição de frequência do gasto de campanha (escala logarítmica,
           candidatos com Gasto $>$ 0). Brasil, 2014.}
  \label{fig:hist_gasto}
\end{figure}

A distribuição é fortemente assimétrica à direita em todos os recortes
analisados. A grande maioria dos candidatos concentra-se em gastos inferiores
a R\$\,50.000, enquanto um pequeno número ultrapassou R\$\,1 milhão. Nenhum
candidato eleito registrou gasto abaixo de R\$\,38.810, sugerindo a existência
de um limiar mínimo de visibilidade para o sucesso eleitoral.

\subsection{Distribuição das Variáveis Qualitativas}

\subsubsection{Gasto por Partido}

As Figuras~\ref{fig:gasto_total_partido} e~\ref{fig:gasto_medio_partido}
mostram, respectivamente, o gasto total e o gasto médio por partido em nível
nacional.

\begin{figure}[t]
  \centering
  \includegraphics[width=\columnwidth]{gasto_total_partido.png}
  \caption{Gasto total de campanha por partido (top partidos). Brasil, 2014.}
  \label{fig:gasto_total_partido}
\end{figure}

\begin{figure}[t]
  \centering
  \includegraphics[width=\columnwidth]{gasto_medio_partido.png}
  \caption{Gasto médio de campanha por partido. Brasil, 2014.}
  \label{fig:gasto_medio_partido}
\end{figure}

A análise por partido revela concentração de recursos em partidos de maior
porte, com significativa variação no gasto médio por candidato. Partidos com
maior número de candidatos tendem a dominar o gasto total, mas não
necessariamente o gasto médio por candidato.

\subsubsection{Gasto por Estado}

As Figuras~\ref{fig:gasto_total_estado} e~\ref{fig:gasto_medio_estado}
apresentam o gasto total e médio por estado.

\begin{figure}[t]
  \centering
  \includegraphics[width=\columnwidth]{gasto_total_estado.png}
  \caption{Gasto total de campanha por estado. Brasil, 2014.}
  \label{fig:gasto_total_estado}
\end{figure}

\begin{figure}[t]
  \centering
  \includegraphics[width=\columnwidth]{gasto_medio_estado.png}
  \caption{Gasto médio de campanha por estado. Brasil, 2014.}
  \label{fig:gasto_medio_estado}
\end{figure}

São Paulo (SP) lidera o gasto total absoluto, reflexo de sua bancada de 70
deputados e consequente maior número de candidatos (1.072). Contudo, no ranking
de gasto médio por candidato, estados menores como Goiás e Amazonas superam SP
e RJ, pois disputam o mesmo número de vagas com um número menor de candidatos,
concentrando mais recursos por candidatura.

\subsection{Correlação entre Gasto e Votos}

\subsubsection{Nível Nacional}

A Figura~\ref{fig:scatter_br} apresenta o diagrama de dispersão entre gasto
de campanha e votos obtidos em nível nacional, com pontos coloridos por
resultado eleitoral. A Figura~\ref{fig:matriz_br} exibe a matriz de correlação
de Pearson.

\begin{figure}[t]
  \centering
  \includegraphics[width=\columnwidth]{gasto_vs_votos.png}
  \caption{Dispersão: Gasto $\times$ Votos, por resultado eleitoral.
           $r = 0{,}593$. Brasil, 2014.}
  \label{fig:scatter_br}
\end{figure}

\begin{figure}[t]
  \centering
  \includegraphics[width=\columnwidth]{br_matriz_correlacao.png}
  \caption{Matriz de correlação de Pearson (Gasto $\times$ Votos). Brasil, 2014.}
  \label{fig:matriz_br}
\end{figure}

O coeficiente de Pearson nacional é $r = 0{,}593$ ($p < 0{,}001$), indicando
correlação positiva moderada a forte. A Tabela~\ref{tab:pearson_br} detalha
os resultados por grupo eleitoral.

\begin{figure}[t]
  \centering
  \tablePlaceholder{Gasto e Votos por resultado eleitoral, com correlação de Pearson --- Brasil.}
  \caption{Gasto e votos por resultado eleitoral, e coeficiente de correlação
           de Pearson entre Gasto e Votos. Brasil, 2014.}
  \label{tab:pearson_br}
\end{figure}

Os candidatos eleitos apresentaram gasto médio de R\$\,1.475.261 (mediana:
R\$\,1.072.042), contra R\$\,119.472 dos suplentes e R\$\,23.316 dos não
eleitos. A diferença de medianas entre eleitos e não eleitos é de $\approx
310\times$, ilustrando que o gasto mediano dos eleitos supera o de quase todos
os candidatos não eleitos.

\subsubsection{Divergências entre Fontes por Estado}

A Figura~\ref{fig:divergentes} exibe a distribuição das 264 divergências de
gasto encontradas entre a base do professor e os dados oficiais do TSE,
estratificadas por estado.

\begin{figure}[t]
  \centering
  \includegraphics[width=\columnwidth]{divergentes_por_estado.png}
  \caption{Divergências de gasto entre dataset do professor e TSE, por estado.
           Brasil, 2014.}
  \label{fig:divergentes}
\end{figure}

A concentração de $\approx 40\%$ das divergências no Rio de Janeiro (107 casos)
é consistente com a hipótese de que os dados do professor foram compilados a
partir de uma versão anterior da base do TSE, possivelmente antes de
retificações realizadas pelos candidatos fluminenses.

\subsection{Análise por Estado: RJ, GO e AM}

\subsubsection{Rio de Janeiro (RJ)}

\begin{figure}[t]
  \centering
  \tablePlaceholder{Estatísticas descritivas de Gasto e Votos --- Rio de Janeiro (RJ).}
  \caption{Estatísticas descritivas das variáveis quantitativas --- RJ, 2014.}
  \label{tab:descritivas_rj}
\end{figure}

O Rio de Janeiro contou com 1.001 candidatos (699 com votos registrados) e 46
eleitos. A mediana de gastos (R\$\,3.324) é 22{,}7$\times$ menor que a média
(R\$\,110.674), assimetria análoga ao padrão nacional. A Figura~\ref{fig:scatter_rj}
e a Tabela~\ref{tab:pearson_rj} apresentam a correlação estadual.

\begin{figure}[t]
  \centering
  \includegraphics[width=\columnwidth]{gasto_vs_votos_RJ.png}
  \caption{Dispersão: Gasto $\times$ Votos --- RJ. $r = 0{,}556$. 2014.}
  \label{fig:scatter_rj}
\end{figure}

\begin{figure}[t]
  \centering
  \tablePlaceholder{Gasto e Votos por resultado eleitoral, com correlação de Pearson --- RJ.}
  \caption{Gasto e votos por resultado eleitoral, e coeficiente de correlação
           de Pearson entre Gasto e Votos. RJ, 2014.}
  \label{tab:pearson_rj}
\end{figure}

A correlação no RJ ($r = 0{,}556$, $p < 0{,}001$) é a mais baixa entre os
três estados analisados, reflexo da heterogeneidade inerente ao estado mais
populoso da amostra (701 candidatos com votos). Os eleitos gastaram em média
R\$\,1.618.134, contra R\$\,17.496 dos não eleitos.

\subsubsection{Goiás (GO)}

\begin{figure}[t]
  \centering
  \tablePlaceholder{Estatísticas descritivas de Gasto e Votos --- Goiás (GO).}
  \caption{Estatísticas descritivas das variáveis quantitativas --- GO, 2014.}
  \label{tab:descritivas_go}
\end{figure}

Goiás apresentou 144 candidatos (77 com votos), com 17 eleitos. A razão
média/mediana é de 125{,}7$\times$ --- a mais extrema entre os estados analisados
--- indicando que poucos candidatos de alto gasto dominam completamente a
distribuição. A Figura~\ref{fig:scatter_go} e a Tabela~\ref{tab:pearson_go}
apresentam a correlação estadual.

\begin{figure}[t]
  \centering
  \includegraphics[width=\columnwidth]{gasto_vs_votos_GO.png}
  \caption{Dispersão: Gasto $\times$ Votos --- GO. $r = 0{,}605$. 2014.}
  \label{fig:scatter_go}
\end{figure}

\begin{figure}[t]
  \centering
  \tablePlaceholder{Gasto e Votos por resultado eleitoral, com correlação de Pearson --- GO.}
  \caption{Gasto e votos por resultado eleitoral, e coeficiente de correlação
           de Pearson entre Gasto e Votos. GO, 2014.}
  \label{tab:pearson_go}
\end{figure}

Os eleitos em GO gastaram em média R\$\,2.517.461 (mediana: R\$\,2.222.030),
contra R\$\,16.712 dos não eleitos --- uma diferença de 151$\times$ nas médias
e de 519$\times$ nas medianas, a mais pronunciada entre os três estados.
A correlação de $r = 0{,}605$ ($p < 0{,}001$) é superior à do RJ.

\subsubsection{Amazonas (AM)}

\begin{figure}[t]
  \centering
  \tablePlaceholder{Estatísticas descritivas de Gasto e Votos --- Amazonas (AM).}
  \caption{Estatísticas descritivas das variáveis quantitativas --- AM, 2014.}
  \label{tab:descritivas_am}
\end{figure}

O Amazonas contou com 80 candidatos (64 com votos), com 8 eleitos. É o estado
com a correlação mais forte entre gasto e votos de toda a análise. A
Figura~\ref{fig:scatter_am} e a Tabela~\ref{tab:pearson_am} apresentam os
resultados.

\begin{figure}[t]
  \centering
  \includegraphics[width=\columnwidth]{gasto_vs_votos_AM.png}
  \caption{Dispersão: Gasto $\times$ Votos --- AM. $r = 0{,}902$. 2014.}
  \label{fig:scatter_am}
\end{figure}

\begin{figure}[t]
  \centering
  \tablePlaceholder{Gasto e Votos por resultado eleitoral, com correlação de Pearson --- AM.}
  \caption{Gasto e votos por resultado eleitoral, e coeficiente de correlação
           de Pearson entre Gasto e Votos. AM, 2014.}
  \label{tab:pearson_am}
\end{figure}

O AM exibiu a correlação mais alta entre os 27 estados ($r = 0{,}902$,
$p < 0{,}001$), sugerindo que o gasto de campanha é um preditor quase linear
dos votos obtidos nesse estado. A diferença de medianas entre eleitos
(R\$\,1.449.783) e não eleitos (R\$\,1.641) é de $\approx 883\times$.

A Figura~\ref{fig:dist_gasto_estados} apresenta as distribuições de gasto e
a Figura~\ref{fig:dist_votos_estados} as distribuições de votos para cada
estado, permitindo comparação visual direta.

\begin{figure*}[t]
  \centering
  \begin{subfigure}[b]{0.32\textwidth}
    \includegraphics[width=\textwidth]{rj_distribuicao_gasto.png}
    \caption{RJ}
  \end{subfigure}
  \hfill
  \begin{subfigure}[b]{0.32\textwidth}
    \includegraphics[width=\textwidth]{go_distribuicao_gasto.png}
    \caption{GO}
  \end{subfigure}
  \hfill
  \begin{subfigure}[b]{0.32\textwidth}
    \includegraphics[width=\textwidth]{am_distribuicao_gasto.png}
    \caption{AM}
  \end{subfigure}
  \caption{Distribuição de frequência do gasto de campanha por estado
           (escala logarítmica). 2014.}
  \label{fig:dist_gasto_estados}
\end{figure*}

\begin{figure*}[t]
  \centering
  \begin{subfigure}[b]{0.32\textwidth}
    \includegraphics[width=\textwidth]{rj_distribuicao_votos.png}
    \caption{RJ}
  \end{subfigure}
  \hfill
  \begin{subfigure}[b]{0.32\textwidth}
    \includegraphics[width=\textwidth]{go_distribuicao_votos.png}
    \caption{GO}
  \end{subfigure}
  \hfill
  \begin{subfigure}[b]{0.32\textwidth}
    \includegraphics[width=\textwidth]{am_distribuicao_votos.png}
    \caption{AM}
  \end{subfigure}
  \caption{Distribuição de frequência de votos obtidos por estado
           (escala logarítmica). 2014.}
  \label{fig:dist_votos_estados}
\end{figure*}

\subsubsection{Síntese da Correlação por Estado}

A Tabela~\ref{tab:pearson_estados} resume os coeficientes de Pearson para todos
os estados, evidenciando a universalidade da correlação positiva entre gasto
e votos.

\begin{figure}[t]
  \centering
  \tablePlaceholder{Coeficiente de Pearson entre Gasto e Votos por estado (todos os 27 estados).}
  \caption{Coeficiente de correlação de Pearson entre Gasto e Votos por estado.
           Todos os estados brasileiros, 2014.}
  \label{tab:pearson_estados}
\end{figure}

A correlação foi positiva em todos os 27 estados ($p < 0{,}001$), variando de
$r = 0{,}902$ no Amazonas a $r = 0{,}475$ em São Paulo. Estados menores tendem
a apresentar correlações mais fortes, possivelmente porque o menor número de
candidatos torna a relação entre gasto e votos mais linear e menos ruidosa.

\subsection{Análise de Fornecedores e Setores Econômicos}

As Figuras~\ref{fig:fornecedores_top10} e~\ref{fig:fornecedores_setores}
apresentam os dez maiores fornecedores e os principais setores econômicos
receptores de gastos de campanha em nível nacional.

\begin{figure}[t]
  \centering
  \includegraphics[width=\columnwidth]{fornecedores_top10.png}
  \caption{Top 10 fornecedores por valor total recebido. Brasil, 2014.}
  \label{fig:fornecedores_top10}
\end{figure}

\begin{figure}[t]
  \centering
  \includegraphics[width=\columnwidth]{fornecedores_setores.png}
  \caption{Distribuição do gasto por setor econômico dos fornecedores. Brasil, 2014.}
  \label{fig:fornecedores_setores}
\end{figure}

A categoria \textbf{\#NULO} (despesas sem CNPJ de fornecedor classificado)
representa 40{,}8\% dos gastos nacionais, 28{,}8\% no RJ e 61{,}3\% em GO,
o que limita a análise setorial. Entre os setores identificados, destacam-se
serviços gráficos e comunicação visual, comitês partidários e serviços de
transporte. O maior fornecedor individual foi a COOPGUANABARA, com
R\$\,3.697.968 em 61 registros --- cooperativa fundada no mesmo ano da
eleição, o que levanta questões sobre sua origem e relação com os candidatos.

No Amazonas, a presença de serviços de \textit{táxi aéreo} entre os principais
setores é uma característica única, refletindo as restrições logísticas de um
estado sem acesso terrestre a grande parte de seu território.

As Figuras~\ref{fig:fornecedores_estados} apresentam os top 10 fornecedores
por estado para comparação regional.

\begin{figure*}[t]
  \centering
  \begin{subfigure}[b]{0.32\textwidth}
    \includegraphics[width=\textwidth]{rj_fornecedores_top10.png}
    \caption{RJ}
  \end{subfigure}
  \hfill
  \begin{subfigure}[b]{0.32\textwidth}
    \includegraphics[width=\textwidth]{go_fornecedores_top10.png}
    \caption{GO}
  \end{subfigure}
  \hfill
  \begin{subfigure}[b]{0.32\textwidth}
    \includegraphics[width=\textwidth]{am_fornecedores_top10.png}
    \caption{AM}
  \end{subfigure}
  \caption{Top 10 fornecedores por valor total recebido, por estado. 2014.}
  \label{fig:fornecedores_estados}
\end{figure*}

\subsection{Análise de Gênero}
\label{subsec:genero}

\subsubsection{Representatividade e Gasto}

As Figuras~\ref{fig:genero_distribuicao}, \ref{fig:genero_representatividade}
e~\ref{fig:genero_taxa} apresentam, respectivamente, a distribuição de
candidatos por gênero, a razão de representatividade financeira e a taxa de
eleição por gênero em nível nacional.

\begin{figure}[t]
  \centering
  \includegraphics[width=\columnwidth]{genero_distribuicao.png}
  \caption{Distribuição de candidatos por gênero: contagem e proporção.
           Brasil, 2014.}
  \label{fig:genero_distribuicao}
\end{figure}

\begin{figure}[t]
  \centering
  \includegraphics[width=\columnwidth]{genero_representatividade.png}
  \caption{Razão de representatividade financeira por gênero (proporção do
           gasto / proporção de candidatos). Brasil, 2014.}
  \label{fig:genero_representatividade}
\end{figure}

\begin{figure}[t]
  \centering
  \includegraphics[width=\columnwidth]{genero_taxa_eleicao.png}
  \caption{Taxa de eleição por gênero. Brasil, 2014.}
  \label{fig:genero_taxa}
\end{figure}

Candidatas mulheres representaram 29{,}6\% do total, mas concentraram apenas
9{,}6\% dos gastos totais de campanha --- uma razão de representatividade de
$0{,}32\times$. Em contrapartida, candidatos homens (70{,}4\% dos candidatos)
concentraram 90{,}4\% dos gastos (razão $1{,}28\times$). Essa assimetria
financeira se reflete nos resultados eleitorais: mulheres corresponderam a
apenas 8{,}0\% dos eleitos, abaixo de sua já reduzida proporção na corrida
eleitoral.

A Figura~\ref{fig:genero_boxplot} apresenta os boxplots de gasto por gênero,
e as Figuras~\ref{fig:genero_estados} mostram as distribuições estaduais.

\begin{figure}[t]
  \centering
  \includegraphics[width=\columnwidth]{genero_boxplot.png}
  \caption{Boxplot do gasto de campanha por gênero (escala logarítmica).
           Brasil, 2014.}
  \label{fig:genero_boxplot}
\end{figure}

\begin{figure*}[t]
  \centering
  \begin{subfigure}[b]{0.32\textwidth}
    \includegraphics[width=\textwidth]{rj_genero_representatividade.png}
    \caption{RJ}
  \end{subfigure}
  \hfill
  \begin{subfigure}[b]{0.32\textwidth}
    \includegraphics[width=\textwidth]{go_genero_representatividade.png}
    \caption{GO}
  \end{subfigure}
  \hfill
  \begin{subfigure}[b]{0.32\textwidth}
    \includegraphics[width=\textwidth]{am_genero_representatividade.png}
    \caption{AM}
  \end{subfigure}
  \caption{Razão de representatividade financeira por gênero, por estado. 2014.}
  \label{fig:genero_estados}
\end{figure*}

A disparidade de gênero se aprofunda nos estados menores: em Goiás, a razão
feminina cai para $0{,}40\times$; no Amazonas, para $0{,}20\times$, indicando
que candidatas amazonenses acessam apenas um quinto dos recursos financeiros
proporcionais ao seu número de candidaturas.

\subsection{Análise de Raça}
\label{subsec:raca}

\subsubsection{Representatividade e Gasto}

As Figuras~\ref{fig:raca_distribuicao}, \ref{fig:raca_representatividade}
e~\ref{fig:raca_taxa} apresentam a distribuição por raça, a razão de
representatividade e a taxa de eleição em nível nacional.

\begin{figure}[t]
  \centering
  \includegraphics[width=\columnwidth]{raca_distribuicao.png}
  \caption{Distribuição de candidatos por raça: contagem e proporção.
           Brasil, 2014.}
  \label{fig:raca_distribuicao}
\end{figure}

\begin{figure}[t]
  \centering
  \includegraphics[width=\columnwidth]{raca_representatividade.png}
  \caption{Razão de representatividade financeira por raça (proporção do
           gasto / proporção de candidatos). Brasil, 2014.}
  \label{fig:raca_representatividade}
\end{figure}

\begin{figure}[t]
  \centering
  \includegraphics[width=\columnwidth]{raca_taxa_eleicao.png}
  \caption{Taxa de eleição por raça. Brasil, 2014.}
  \label{fig:raca_taxa}
\end{figure}

Candidatos brancos são o único grupo com razão de representatividade superior
a 1 ($1{,}40\times$), acessando 40\% mais recursos do que sua proporção de
candidatos justificaria. Todos os demais grupos recebem proporcionalmente menos:
pardos ($0{,}46\times$), pretos ($0{,}38\times$), amarelos e indígenas
($0{,}05\times$). Candidatos indígenas são virtualmente ausentes do
financiamento eleitoral.

A Figura~\ref{fig:raca_boxplot} apresenta os boxplots de gasto por raça, e as
Figuras~\ref{fig:raca_estados} mostram as razões estaduais.

\begin{figure}[t]
  \centering
  \includegraphics[width=\columnwidth]{raca_boxplot.png}
  \caption{Boxplot do gasto de campanha por raça (escala logarítmica).
           Brasil, 2014.}
  \label{fig:raca_boxplot}
\end{figure}

\begin{figure*}[t]
  \centering
  \begin{subfigure}[b]{0.32\textwidth}
    \includegraphics[width=\textwidth]{rj_raca_representatividade.png}
    \caption{RJ}
  \end{subfigure}
  \hfill
  \begin{subfigure}[b]{0.32\textwidth}
    \includegraphics[width=\textwidth]{go_raca_representatividade.png}
    \caption{GO}
  \end{subfigure}
  \hfill
  \begin{subfigure}[b]{0.32\textwidth}
    \includegraphics[width=\textwidth]{am_raca_representatividade.png}
    \caption{AM}
  \end{subfigure}
  \caption{Razão de representatividade financeira por raça, por estado. 2014.}
  \label{fig:raca_estados}
\end{figure*}

Destaca-se a situação dos candidatos negros em Goiás: com 5{,}6\% dos
candidatos, concentraram apenas 0{,}1\% dos gastos (razão $0{,}02\times$) e
nenhum foi eleito. No Amazonas --- o estado com a maior proporção de população
indígena do país --- apenas um candidato indígena foi registrado, sem votos
expressivos, o que torna essa ausência ainda mais notável.

% =========================================================================
\section{Conclusão}

Este trabalho realizou uma análise estatística descritiva abrangente do dataset
de eleições brasileiras de 2014, investigando a relação entre gasto de campanha,
votos obtidos e resultado eleitoral, com análises complementares de fornecedores,
gênero e raça.

A distribuição de gastos é fortemente assimétrica à direita em todos os recortes
analisados. Em nível nacional, a mediana (R\$\,5.804) é 22$\times$ inferior à
média (R\$\,184.664), tornando a mediana o indicador mais adequado do candidato
típico. Os desvios-padrão superam as médias em todos os estados, confirmando a
inadequação da média como medida única de posição. A assimetria é ainda mais
extrema em Goiás (razão de 125,7$\times$) e no Amazonas (42,7$\times$).

A correlação positiva entre gasto de campanha e votos obtidos foi verificada
em todos os 27 estados brasileiros ($p < 0{,}001$), variando de $r = 0{,}902$
no Amazonas a $r = 0{,}475$ em São Paulo. A relação não é determinística ---
altos gastos não garantem eleição, mas a ausência de gasto parece efetivamente
impedir o sucesso eleitoral: nenhum candidato eleito registrou gasto abaixo de
R\$\,38.810. Estados menores tendem a apresentar correlações mais fortes,
possivelmente porque a disputa por vagas é mais concentrada e a visibilidade
financeira tem impacto mais direto no resultado.

A análise de equidade revela disparidades estruturais no acesso ao financiamento
eleitoral. No recorte de gênero, candidatas mulheres concentraram 9{,}6\% dos
gastos sendo 29{,}6\% dos candidatos (razão $0{,}32\times$) e elegeram apenas
8{,}0\% dos eleitos. No recorte racial, candidatos brancos foram o único grupo
com acesso proporcionalmente superior ao seu tamanho (razão $1{,}40\times$),
enquanto candidatos indígenas acessaram apenas 5\% do esperado. Esses resultados
sugerem que o financiamento de campanha atua como barreira estrutural à
representação de grupos sub-representados no Legislativo Federal.

Os dados analisados referem-se exclusivamente às eleições de 2014 para Deputado
Federal. Generalizações para outros cargos ou anos eleitorais requerem análises
específicas, dado que as dinâmicas de financiamento e representatividade podem
variar significativamente conforme o eleitorado e as regras eleitorais vigentes.

\end{document}
